
\documentclass{article}
\usepackage{colortbl}
\usepackage{makecell}
\usepackage{multirow}
\usepackage{supertabular}

\begin{document}

\newcounter{utterance}

\centering \large Interaction Transcript for game `matchit\_ascii', experiment `same\_grid', episode 3 with deepseek{-}v4{-}pro.
\vspace{24pt}

{ \footnotesize  \setcounter{utterance}{1}
\setlength{\tabcolsep}{0pt}
\begin{supertabular}{c@{$\;$}|p{.15\linewidth}@{}p{.15\linewidth}p{.15\linewidth}p{.15\linewidth}p{.15\linewidth}p{.15\linewidth}}
    \# & $\;$A & \multicolumn{4}{c}{Game Master} & $\;\:$B\\
    \hline

    \theutterance \stepcounter{utterance}  
    & & \multicolumn{4}{p{0.6\linewidth}}{
        \cellcolor[rgb]{0.9,0.9,0.9}{
            \makecell[{{p{\linewidth}}}]{
                \texttt{\tiny{[P1$\langle$GM]}}
                \texttt{Zúčastňujete sa spoločnej hry na hádanie.} \\
\texttt{Cieľom je zistiť, či je táto mriežka a iný obrázok, ktorý vidím len ja, rovnaké.} \\
\texttt{Najskôr opíšte svoju mriežku. Potom poskytnem svoj opis a môžeme si navzájom klásť otázky o mriežkach, aby sme zistili, či sú rovnaké. Teraz začnite svoj krátky opis mriežky s "POPIS:", za ktorým nasleduje opis. Nepridávajte nič iné. Toto je vaša mriežka:} \\
\texttt{X X X X X} \\
\texttt{▢ ▢ ▢ ▢ X} \\
\texttt{▢ ▢ X X X} \\
\texttt{▢ ▢ ▢ ▢ X} \\
\texttt{▢ ▢ ▢ ▢ X} \\
            }
        }
    }
    & \\ \\

    \theutterance \stepcounter{utterance}  
    & \multicolumn{4}{p{0.6\linewidth}}{
        \cellcolor[rgb]{0.9,0.9,0.9}{
            \makecell[{{p{\linewidth}}}]{
                \texttt{\tiny{[P1$\rangle$GM]}}
                \texttt{POPIS: Mriežka má 5 riadkov a 5 stĺpcov. Prvý riadok je plný X. V druhom riadku sú prvé štyri bunky prázdne (▢) a piata bunka je X. V treťom riadku sú prvé dve bunky prázdne (▢) a zvyšné tri bunky sú X. Vo štvrtom riadku sú prvé štyri bunky prázdne (▢) a piata bunka je X. V piatom riadku sú prvé štyri bunky prázdne (▢) a piata bunka je X.} \\
            }
        }
    }
    & & \\ \\

    \theutterance \stepcounter{utterance}  
    & & & \multicolumn{2}{p{0.3\linewidth}}{
        \cellcolor[rgb]{0.9,0.9,0.9}{
            \makecell[{{p{\linewidth}}}]{
                \texttt{\tiny{[GM$|$GM]}}
                \texttt{continue} \\
            }
        }
    }
    & & \\ \\

    \theutterance \stepcounter{utterance}  
    & & \multicolumn{4}{p{0.6\linewidth}}{
        \cellcolor[rgb]{0.9,0.9,0.9}{
            \makecell[{{p{\linewidth}}}]{
                \texttt{\tiny{[GM$\rangle$P2]}}
                \texttt{Zúčastňujete sa spoločnej hry na hádanie.} \\
\texttt{Cieľom je zistiť, či je táto mriežka a iný obrázok, ktorý vidím len ja, rovnaké.} \\
\texttt{Najskôr opíšte svoju mriežku. Potom poskytnem svoj opis a môžeme si navzájom klásť otázky o mriežkach, aby sme zistili, či sú rovnaké. Teraz začnite svoj krátky opis mriežky s "POPIS:", za ktorým nasleduje opis. Nepridávajte nič iné. Toto je vaša mriežka:} \\
\texttt{X X X X X} \\
\texttt{▢ ▢ ▢ ▢ X} \\
\texttt{▢ ▢ X X X} \\
\texttt{▢ ▢ ▢ ▢ X} \\
\texttt{▢ ▢ ▢ ▢ X} \\
            }
        }
    }
    & \\ \\

    \theutterance \stepcounter{utterance}  
    & & & \multicolumn{4}{p{0.6\linewidth}}{
        \cellcolor[rgb]{0.9,0.9,0.9}{
            \makecell[{{p{\linewidth}}}]{
                \texttt{\tiny{[GM$\langle$P2]}}
                \texttt{POPIS: Mriežka má 5 riadkov a 5 stĺpcov. Prvý riadok je celý vyplnený X. V druhom riadku sú prvé štyri políčka prázdne (▢) a piate políčko je X. V treťom riadku sú prvé dve políčka prázdne a zvyšné tri sú X. Vo štvrtom riadku sú prvé štyri políčka prázdne a piate políčko je X. V piatom riadku sú prvé štyri políčka prázdne a piate políčko je X.} \\
            }
        }
    }
     \\ \\

    \theutterance \stepcounter{utterance}  
    & & & \multicolumn{2}{p{0.3\linewidth}}{
        \cellcolor[rgb]{0.9,0.9,0.9}{
            \makecell[{{p{\linewidth}}}]{
                \texttt{\tiny{[GM$|$GM]}}
                \texttt{continue} \\
            }
        }
    }
    & & \\ \\

    \theutterance \stepcounter{utterance}  
    & & \multicolumn{4}{p{0.6\linewidth}}{
        \cellcolor[rgb]{0.9,0.9,0.9}{
            \makecell[{{p{\linewidth}}}]{
                \texttt{\tiny{[GM$\rangle$P2]}}
                \texttt{Toto je môj POPISPOPIS: Mriežka má 5 riadkov a 5 stĺpcov. Prvý riadok je plný X. V druhom riadku sú prvé štyri bunky prázdne (▢) a piata bunka je X. V treťom riadku sú prvé dve bunky prázdne (▢) a zvyšné tri bunky sú X. Vo štvrtom riadku sú prvé štyri bunky prázdne (▢) a piata bunka je X. V piatom riadku sú prvé štyri bunky prázdne (▢) a piata bunka je X.} \\
\texttt{Teraz položte otázku, aby ste zistili nové aspekty mojej mriežky, ktoré sa môžu líšiť od vašej mriežky. Začnite s "OTÁZKA:" a nepridávajte nič iné.} \\
            }
        }
    }
    & \\ \\

    \theutterance \stepcounter{utterance}  
    & & & \multicolumn{4}{p{0.6\linewidth}}{
        \cellcolor[rgb]{0.9,0.9,0.9}{
            \makecell[{{p{\linewidth}}}]{
                \texttt{\tiny{[GM$\langle$P2]}}
                \texttt{OTÁZKA: Sú niektoré z prázdnych políčok (▢) v tvojej mriežke vyplnené špecifickým symbolom, farbou alebo číslom, ktoré nie je viditeľné v mojom zobrazení, alebo sú všetky prázdne políčka úplne identické?} \\
            }
        }
    }
     \\ \\

    \theutterance \stepcounter{utterance}  
    & & & \multicolumn{2}{p{0.3\linewidth}}{
        \cellcolor[rgb]{0.9,0.9,0.9}{
            \makecell[{{p{\linewidth}}}]{
                \texttt{\tiny{[GM$|$GM]}}
                \texttt{continue} \\
            }
        }
    }
    & & \\ \\

    \theutterance \stepcounter{utterance}  
    & & \multicolumn{4}{p{0.6\linewidth}}{
        \cellcolor[rgb]{0.9,0.9,0.9}{
            \makecell[{{p{\linewidth}}}]{
                \texttt{\tiny{[P1$\langle$GM]}}
                \texttt{Toto je môj POPISPOPIS: Mriežka má 5 riadkov a 5 stĺpcov. Prvý riadok je celý vyplnený X. V druhom riadku sú prvé štyri políčka prázdne (▢) a piate políčko je X. V treťom riadku sú prvé dve políčka prázdne a zvyšné tri sú X. Vo štvrtom riadku sú prvé štyri políčka prázdne a piate políčko je X. V piatom riadku sú prvé štyri políčka prázdne a piate políčko je X.} \\
\texttt{OTÁZKA: Sú niektoré z prázdnych políčok (▢) v tvojej mriežke vyplnené špecifickým symbolom, farbou alebo číslom, ktoré nie je viditeľné v mojom zobrazení, alebo sú všetky prázdne políčka úplne identické?Svoju odpoveď začnite s "ODPOVEĎ:" a nepridávajte nič iné.} \\
            }
        }
    }
    & \\ \\

    \theutterance \stepcounter{utterance}  
    & \multicolumn{4}{p{0.6\linewidth}}{
        \cellcolor[rgb]{0.9,0.9,0.9}{
            \makecell[{{p{\linewidth}}}]{
                \texttt{\tiny{[P1$\rangle$GM]}}
                \texttt{ODPOVEĎ: Všetky prázdne políčka (▢) v mojej mriežke sú úplne identické a neobsahujú žiadne ďalšie symboly, farby ani čísla.} \\
            }
        }
    }
    & & \\ \\

    \theutterance \stepcounter{utterance}  
    & & & \multicolumn{2}{p{0.3\linewidth}}{
        \cellcolor[rgb]{0.9,0.9,0.9}{
            \makecell[{{p{\linewidth}}}]{
                \texttt{\tiny{[GM$|$GM]}}
                \texttt{continue} \\
            }
        }
    }
    & & \\ \\

    \theutterance \stepcounter{utterance}  
    & & \multicolumn{4}{p{0.6\linewidth}}{
        \cellcolor[rgb]{0.9,0.9,0.9}{
            \makecell[{{p{\linewidth}}}]{
                \texttt{\tiny{[P1$\langle$GM]}}
                \texttt{Teraz položte otázku, aby ste zistili nové aspekty mojej mriežky, ktoré sa môžu líšiť od vašej mriežky. Začnite s "OTÁZKA:" a nepridávajte nič iné.} \\
            }
        }
    }
    & \\ \\

    \theutterance \stepcounter{utterance}  
    & \multicolumn{4}{p{0.6\linewidth}}{
        \cellcolor[rgb]{0.9,0.9,0.9}{
            \makecell[{{p{\linewidth}}}]{
                \texttt{\tiny{[P1$\rangle$GM]}}
                \texttt{OTÁZKA: Sú všetky X v tvojej mriežke rovnakej farby, veľkosti a štýlu, alebo sa niektoré z nich odlišujú (napríklad sú niektoré tučné, inej farby, alebo obsahujú dodatočné prvky)?} \\
            }
        }
    }
    & & \\ \\

    \theutterance \stepcounter{utterance}  
    & & & \multicolumn{2}{p{0.3\linewidth}}{
        \cellcolor[rgb]{0.9,0.9,0.9}{
            \makecell[{{p{\linewidth}}}]{
                \texttt{\tiny{[GM$|$GM]}}
                \texttt{continue} \\
            }
        }
    }
    & & \\ \\

    \theutterance \stepcounter{utterance}  
    & & \multicolumn{4}{p{0.6\linewidth}}{
        \cellcolor[rgb]{0.9,0.9,0.9}{
            \makecell[{{p{\linewidth}}}]{
                \texttt{\tiny{[GM$\rangle$P2]}}
                \texttt{ODPOVEĎ: Všetky prázdne políčka (▢) v mojej mriežke sú úplne identické a neobsahujú žiadne ďalšie symboly, farby ani čísla.} \\
\texttt{OTÁZKA: Sú všetky X v tvojej mriežke rovnakej farby, veľkosti a štýlu, alebo sa niektoré z nich odlišujú (napríklad sú niektoré tučné, inej farby, alebo obsahujú dodatočné prvky)?Svoju odpoveď začnite s "ODPOVEĎ:" a nepridávajte nič iné.} \\
            }
        }
    }
    & \\ \\

    \theutterance \stepcounter{utterance}  
    & & & \multicolumn{4}{p{0.6\linewidth}}{
        \cellcolor[rgb]{0.9,0.9,0.9}{
            \makecell[{{p{\linewidth}}}]{
                \texttt{\tiny{[GM$\langle$P2]}}
                \texttt{ODPOVEĎ: Všetky X v mojej mriežke sú úplne rovnaké – majú rovnakú farbu, veľkosť a štýl, bez akýchkoľvek dodatočných prvkov.} \\
            }
        }
    }
     \\ \\

    \theutterance \stepcounter{utterance}  
    & & & \multicolumn{2}{p{0.3\linewidth}}{
        \cellcolor[rgb]{0.9,0.9,0.9}{
            \makecell[{{p{\linewidth}}}]{
                \texttt{\tiny{[GM$|$GM]}}
                \texttt{continue} \\
            }
        }
    }
    & & \\ \\

    \theutterance \stepcounter{utterance}  
    & & \multicolumn{4}{p{0.6\linewidth}}{
        \cellcolor[rgb]{0.9,0.9,0.9}{
            \makecell[{{p{\linewidth}}}]{
                \texttt{\tiny{[GM$\rangle$P2]}}
                \texttt{Teraz položte otázku, aby ste zistili nové aspekty mojej mriežky, ktoré sa môžu líšiť od vašej mriežky. Začnite s "OTÁZKA:" a nepridávajte nič iné.} \\
            }
        }
    }
    & \\ \\

    \theutterance \stepcounter{utterance}  
    & & & \multicolumn{4}{p{0.6\linewidth}}{
        \cellcolor[rgb]{0.9,0.9,0.9}{
            \makecell[{{p{\linewidth}}}]{
                \texttt{\tiny{[GM$\langle$P2]}}
                \texttt{OTÁZKA: Je tvoja mriežka niečím ohraničená alebo má vyznačené okraje špecifickým spôsobom (napr. hrubšou čiarou, dvojitou čiarou, farebným rámčekom), ktorý by sa mohol líšiť od štandardnej čiernej mriežky?} \\
            }
        }
    }
     \\ \\

    \theutterance \stepcounter{utterance}  
    & & & \multicolumn{2}{p{0.3\linewidth}}{
        \cellcolor[rgb]{0.9,0.9,0.9}{
            \makecell[{{p{\linewidth}}}]{
                \texttt{\tiny{[GM$|$GM]}}
                \texttt{continue} \\
            }
        }
    }
    & & \\ \\

    \theutterance \stepcounter{utterance}  
    & & \multicolumn{4}{p{0.6\linewidth}}{
        \cellcolor[rgb]{0.9,0.9,0.9}{
            \makecell[{{p{\linewidth}}}]{
                \texttt{\tiny{[P1$\langle$GM]}}
                \texttt{ODPOVEĎ: Všetky X v mojej mriežke sú úplne rovnaké – majú rovnakú farbu, veľkosť a štýl, bez akýchkoľvek dodatočných prvkov.} \\
\texttt{OTÁZKA: Je tvoja mriežka niečím ohraničená alebo má vyznačené okraje špecifickým spôsobom (napr. hrubšou čiarou, dvojitou čiarou, farebným rámčekom), ktorý by sa mohol líšiť od štandardnej čiernej mriežky?Svoju odpoveď začnite s "ODPOVEĎ:" a nepridávajte nič iné.} \\
            }
        }
    }
    & \\ \\

    \theutterance \stepcounter{utterance}  
    & \multicolumn{4}{p{0.6\linewidth}}{
        \cellcolor[rgb]{0.9,0.9,0.9}{
            \makecell[{{p{\linewidth}}}]{
                \texttt{\tiny{[P1$\rangle$GM]}}
                \texttt{ODPOVEĎ: Moja mriežka je ohraničená štandardnou jednoduchou čiernou čiarou, bez hrubších, dvojitých alebo farebných rámčekov.} \\
            }
        }
    }
    & & \\ \\

    \theutterance \stepcounter{utterance}  
    & & & \multicolumn{2}{p{0.3\linewidth}}{
        \cellcolor[rgb]{0.9,0.9,0.9}{
            \makecell[{{p{\linewidth}}}]{
                \texttt{\tiny{[GM$|$GM]}}
                \texttt{continue} \\
            }
        }
    }
    & & \\ \\

    \theutterance \stepcounter{utterance}  
    & & \multicolumn{4}{p{0.6\linewidth}}{
        \cellcolor[rgb]{0.9,0.9,0.9}{
            \makecell[{{p{\linewidth}}}]{
                \texttt{\tiny{[P1$\langle$GM]}}
                \texttt{Teraz položte otázku, aby ste zistili nové aspekty mojej mriežky, ktoré sa môžu líšiť od vašej mriežky. Začnite s "OTÁZKA:" a nepridávajte nič iné.} \\
            }
        }
    }
    & \\ \\

    \theutterance \stepcounter{utterance}  
    & \multicolumn{4}{p{0.6\linewidth}}{
        \cellcolor[rgb]{0.9,0.9,0.9}{
            \makecell[{{p{\linewidth}}}]{
                \texttt{\tiny{[P1$\rangle$GM]}}
                \texttt{OTÁZKA: Je v tvojej mriežke každé políčko dokonale štvorcové a rovnakej veľkosti, alebo sú niektoré políčka obdĺžnikové či inak proporčne odlišné?} \\
            }
        }
    }
    & & \\ \\

    \theutterance \stepcounter{utterance}  
    & & & \multicolumn{2}{p{0.3\linewidth}}{
        \cellcolor[rgb]{0.9,0.9,0.9}{
            \makecell[{{p{\linewidth}}}]{
                \texttt{\tiny{[GM$|$GM]}}
                \texttt{continue} \\
            }
        }
    }
    & & \\ \\

    \theutterance \stepcounter{utterance}  
    & & \multicolumn{4}{p{0.6\linewidth}}{
        \cellcolor[rgb]{0.9,0.9,0.9}{
            \makecell[{{p{\linewidth}}}]{
                \texttt{\tiny{[GM$\rangle$P2]}}
                \texttt{ODPOVEĎ: Moja mriežka je ohraničená štandardnou jednoduchou čiernou čiarou, bez hrubších, dvojitých alebo farebných rámčekov.} \\
\texttt{OTÁZKA: Je v tvojej mriežke každé políčko dokonale štvorcové a rovnakej veľkosti, alebo sú niektoré políčka obdĺžnikové či inak proporčne odlišné?Svoju odpoveď začnite s "ODPOVEĎ:" a nepridávajte nič iné.} \\
            }
        }
    }
    & \\ \\

    \theutterance \stepcounter{utterance}  
    & & & \multicolumn{4}{p{0.6\linewidth}}{
        \cellcolor[rgb]{0.9,0.9,0.9}{
            \makecell[{{p{\linewidth}}}]{
                \texttt{\tiny{[GM$\langle$P2]}}
                \texttt{ODPOVEĎ: V mojej mriežke je každé políčko dokonale štvorcové a všetky políčka majú presne rovnakú veľkosť, bez akýchkoľvek proporčných odlišností.} \\
            }
        }
    }
     \\ \\

    \theutterance \stepcounter{utterance}  
    & & & \multicolumn{2}{p{0.3\linewidth}}{
        \cellcolor[rgb]{0.9,0.9,0.9}{
            \makecell[{{p{\linewidth}}}]{
                \texttt{\tiny{[GM$|$GM]}}
                \texttt{continue} \\
            }
        }
    }
    & & \\ \\

    \theutterance \stepcounter{utterance}  
    & & \multicolumn{4}{p{0.6\linewidth}}{
        \cellcolor[rgb]{0.9,0.9,0.9}{
            \makecell[{{p{\linewidth}}}]{
                \texttt{\tiny{[GM$\rangle$P2]}}
                \texttt{Teraz položte otázku, aby ste zistili nové aspekty mojej mriežky, ktoré sa môžu líšiť od vašej mriežky. Začnite s "OTÁZKA:" a nepridávajte nič iné.} \\
            }
        }
    }
    & \\ \\

    \theutterance \stepcounter{utterance}  
    & & & \multicolumn{4}{p{0.6\linewidth}}{
        \cellcolor[rgb]{0.9,0.9,0.9}{
            \makecell[{{p{\linewidth}}}]{
                \texttt{\tiny{[GM$\langle$P2]}}
                \texttt{OTÁZKA: Je pozadie tvojej mriežky (za prázdnymi políčkami aj za X) úplne biele, alebo má nejaký iný farebný podklad, textúru či vzor?} \\
            }
        }
    }
     \\ \\

    \theutterance \stepcounter{utterance}  
    & & & \multicolumn{2}{p{0.3\linewidth}}{
        \cellcolor[rgb]{0.9,0.9,0.9}{
            \makecell[{{p{\linewidth}}}]{
                \texttt{\tiny{[GM$|$GM]}}
                \texttt{continue} \\
            }
        }
    }
    & & \\ \\

    \theutterance \stepcounter{utterance}  
    & & \multicolumn{4}{p{0.6\linewidth}}{
        \cellcolor[rgb]{0.9,0.9,0.9}{
            \makecell[{{p{\linewidth}}}]{
                \texttt{\tiny{[P1$\langle$GM]}}
                \texttt{ODPOVEĎ: V mojej mriežke je každé políčko dokonale štvorcové a všetky políčka majú presne rovnakú veľkosť, bez akýchkoľvek proporčných odlišností.} \\
\texttt{OTÁZKA: Je pozadie tvojej mriežky (za prázdnymi políčkami aj za X) úplne biele, alebo má nejaký iný farebný podklad, textúru či vzor?Svoju odpoveď začnite s "ODPOVEĎ:" a nepridávajte nič iné.} \\
            }
        }
    }
    & \\ \\

    \theutterance \stepcounter{utterance}  
    & \multicolumn{4}{p{0.6\linewidth}}{
        \cellcolor[rgb]{0.9,0.9,0.9}{
            \makecell[{{p{\linewidth}}}]{
                \texttt{\tiny{[P1$\rangle$GM]}}
                \texttt{ODPOVEĎ: Pozadie mojej mriežky je úplne biele, bez akejkoľvek textúry, vzoru alebo iného farebného podkladu.} \\
            }
        }
    }
    & & \\ \\

    \theutterance \stepcounter{utterance}  
    & & & \multicolumn{2}{p{0.3\linewidth}}{
        \cellcolor[rgb]{0.9,0.9,0.9}{
            \makecell[{{p{\linewidth}}}]{
                \texttt{\tiny{[GM$|$GM]}}
                \texttt{continue} \\
            }
        }
    }
    & & \\ \\

    \theutterance \stepcounter{utterance}  
    & & \multicolumn{4}{p{0.6\linewidth}}{
        \cellcolor[rgb]{0.9,0.9,0.9}{
            \makecell[{{p{\linewidth}}}]{
                \texttt{\tiny{[P1$\langle$GM]}}
                \texttt{Teraz položte otázku, aby ste zistili nové aspekty mojej mriežky, ktoré sa môžu líšiť od vašej mriežky. Začnite s "OTÁZKA:" a nepridávajte nič iné.} \\
            }
        }
    }
    & \\ \\

    \theutterance \stepcounter{utterance}  
    & \multicolumn{4}{p{0.6\linewidth}}{
        \cellcolor[rgb]{0.9,0.9,0.9}{
            \makecell[{{p{\linewidth}}}]{
                \texttt{\tiny{[P1$\rangle$GM]}}
                \texttt{OTÁZKA: Sú v tvojej mriežke medzi jednotlivými políčkami viditeľné čiary, ktoré ich oddeľujú, alebo sú políčka oddelené len implicitne bez vyznačených deliacich čiar?} \\
            }
        }
    }
    & & \\ \\

    \theutterance \stepcounter{utterance}  
    & & & \multicolumn{2}{p{0.3\linewidth}}{
        \cellcolor[rgb]{0.9,0.9,0.9}{
            \makecell[{{p{\linewidth}}}]{
                \texttt{\tiny{[GM$|$GM]}}
                \texttt{continue} \\
            }
        }
    }
    & & \\ \\

    \theutterance \stepcounter{utterance}  
    & & \multicolumn{4}{p{0.6\linewidth}}{
        \cellcolor[rgb]{0.9,0.9,0.9}{
            \makecell[{{p{\linewidth}}}]{
                \texttt{\tiny{[GM$\rangle$P2]}}
                \texttt{ODPOVEĎ: Pozadie mojej mriežky je úplne biele, bez akejkoľvek textúry, vzoru alebo iného farebného podkladu.} \\
\texttt{OTÁZKA: Sú v tvojej mriežke medzi jednotlivými políčkami viditeľné čiary, ktoré ich oddeľujú, alebo sú políčka oddelené len implicitne bez vyznačených deliacich čiar?Svoju odpoveď začnite s "ODPOVEĎ:" a nepridávajte nič iné.} \\
            }
        }
    }
    & \\ \\

    \theutterance \stepcounter{utterance}  
    & & & \multicolumn{4}{p{0.6\linewidth}}{
        \cellcolor[rgb]{0.9,0.9,0.9}{
            \makecell[{{p{\linewidth}}}]{
                \texttt{\tiny{[GM$\langle$P2]}}
                \texttt{ODPOVEĎ: V mojej mriežke sú medzi všetkými jednotlivými políčkami viditeľné tenké čierne deliace čiary, takže každé políčko je explicitne ohraničené.} \\
            }
        }
    }
     \\ \\

    \theutterance \stepcounter{utterance}  
    & & & \multicolumn{2}{p{0.3\linewidth}}{
        \cellcolor[rgb]{0.9,0.9,0.9}{
            \makecell[{{p{\linewidth}}}]{
                \texttt{\tiny{[GM$|$GM]}}
                \texttt{continue} \\
            }
        }
    }
    & & \\ \\

    \theutterance \stepcounter{utterance}  
    & & \multicolumn{4}{p{0.6\linewidth}}{
        \cellcolor[rgb]{0.9,0.9,0.9}{
            \makecell[{{p{\linewidth}}}]{
                \texttt{\tiny{[GM$\rangle$P2]}}
                \texttt{Teraz sa rozhodnite. Čo si myslíte: je vaša mriežka a druhá opísaná mriežka rovnaká? Napíšte "ROZHODNUTIE: rovnaká mriežka" ak si myslíte, že ide o rovnakú mriežku, alebo "ROZHODNUTIE: rôzne mriežky" ak si myslíte, že ide o rôzne mriežky. Nepridávajte nič iné.} \\
            }
        }
    }
    & \\ \\

    \theutterance \stepcounter{utterance}  
    & & & \multicolumn{4}{p{0.6\linewidth}}{
        \cellcolor[rgb]{0.9,0.9,0.9}{
            \makecell[{{p{\linewidth}}}]{
                \texttt{\tiny{[GM$\langle$P2]}}
                \texttt{ROZHODNUTIE: rovnaká mriežka} \\
            }
        }
    }
     \\ \\

    \theutterance \stepcounter{utterance}  
    & & & \multicolumn{2}{p{0.3\linewidth}}{
        \cellcolor[rgb]{0.9,0.9,0.9}{
            \makecell[{{p{\linewidth}}}]{
                \texttt{\tiny{[GM$|$GM]}}
                \texttt{continue} \\
            }
        }
    }
    & & \\ \\

    \theutterance \stepcounter{utterance}  
    & & & \multicolumn{2}{p{0.3\linewidth}}{
        \cellcolor[rgb]{0.9,0.9,0.9}{
            \makecell[{{p{\linewidth}}}]{
                \texttt{\tiny{[GM$|$GM]}}
                \texttt{success} \\
            }
        }
    }
    & & \\ \\

    \theutterance \stepcounter{utterance}  
    & & \multicolumn{4}{p{0.6\linewidth}}{
        \cellcolor[rgb]{0.9,0.9,0.9}{
            \makecell[{{p{\linewidth}}}]{
                \texttt{\tiny{[P1$\langle$GM]}}
                \texttt{ODPOVEĎ: V mojej mriežke sú medzi všetkými jednotlivými políčkami viditeľné tenké čierne deliace čiary, takže každé políčko je explicitne ohraničené.} \\
\texttt{Teraz sa rozhodnite. Čo si myslíte: je vaša mriežka a druhá opísaná mriežka rovnaká? Napíšte "ROZHODNUTIE: rovnaká mriežka" ak si myslíte, že ide o rovnakú mriežku, alebo "ROZHODNUTIE: rôzne mriežky" ak si myslíte, že ide o rôzne mriežky. Nepridávajte nič iné.} \\
            }
        }
    }
    & \\ \\

    \theutterance \stepcounter{utterance}  
    & \multicolumn{4}{p{0.6\linewidth}}{
        \cellcolor[rgb]{0.9,0.9,0.9}{
            \makecell[{{p{\linewidth}}}]{
                \texttt{\tiny{[P1$\rangle$GM]}}
                \texttt{ROZHODNUTIE: rovnaká mriežka} \\
            }
        }
    }
    & & \\ \\

    \theutterance \stepcounter{utterance}  
    & & & \multicolumn{2}{p{0.3\linewidth}}{
        \cellcolor[rgb]{0.9,0.9,0.9}{
            \makecell[{{p{\linewidth}}}]{
                \texttt{\tiny{[GM$|$GM]}}
                \texttt{continue} \\
            }
        }
    }
    & & \\ \\

    \theutterance \stepcounter{utterance}  
    & & & \multicolumn{2}{p{0.3\linewidth}}{
        \cellcolor[rgb]{0.9,0.9,0.9}{
            \makecell[{{p{\linewidth}}}]{
                \texttt{\tiny{[GM$|$GM]}}
                \texttt{success} \\
            }
        }
    }
    & & \\ \\

\end{supertabular}
}

\end{document}
